% !PRINTED 322
\weberpage{322}
%<<CONT
schwinden, wenn (4) zwei gleiche Wurzeln hat, und dies ist nach XIV nur dann
möglich, wenn $f(x,y)$ die $n^{\mathrm{te}}$ Potenz einer linearen Function
ist. In diesem Falle verschwindet die Determinante (5) identisch: Wenn wir
also von diesem Falle absehen, so folgt, \gesperrt{dass für eine Gleichung mit nur
reellen Wurzeln die Gleichung, die man durch Nullsetzen der Hesse'schen
Determinante erhält, nur imaginäre Wurzeln hat.}

Wollen wir zur inhomogenen Darstellung zurückkehren, so können wir mit
Benutzung des Euler'schen Satzes
\weq{6}{\begin{aligned}
xf'(x) + yf'(y) &= nf(x)\\
xf''(x,x) + yf''(x,y) &= (n-1)f'(x)\\
xf''(y,x) + yf''(y,y) &= (n-1)f'(y)
\end{aligned}}
setzen, und damit $f'(y), f''(x,y), f''(y,y)$ eliminiren. Man erhält so,
wenn man $y = 1$ setzt,
\weq{7}{\begin{aligned}
f'(y) &= nf(x) - xf'(x)\\
f''(x,y) &= (n-1)f'(x) - xf''(x)\\
f''(y,y) &= n(n-1)f(x) - 2(n-1)xf'(x) + x^2 f''(x),
\end{aligned}}
also wenn
\weq{8}{(n-1)H(x,y) = f''(x,x)f''(y,y) - f''(x,y)^2}
gesetzt wird:
\weq{9}{H(x,y) = H(x) = nf(x)f''(x) - (n-1)f'(x)^2.}
$H(x)$ ist in Bezug auf $x$ höchstens vom Grade $2n-4$.

\begin{center}
\S.\,107.

Die Sätze von Laguerre für Gleichungen mit nur reellen Wurzeln.
\end{center}

Von den zuletzt abgeleiteten Sätzen hat \gesperrt{Laguerre} eine Anwendung
auf Gleichungen mit nur reellen Wurzeln gemacht, die wir noch kennen lernen
wollen.

Es sei $f(x) = 0$ eine Gleichung $n^{\mathrm{ten}}$ Grades mit $n$ \gesperrt{reellen}
Wurzeln $\alpha, \alpha_1, \alpha_2 \dots$, \gesperrt{die wir von einander verschieden
voraussetzen}.

Ist $x$ eine veränderliche Grösse, so haben wir die schon mehrfach
angewandte Formel [\S.\,14, (11)]:
%CONT>>
